\documentclass[
    reprint,
    amsmath,amssymb,
    aps,
    pra,
    superscriptaddress,
    floatfix
]{revtex4-2}

\usepackage{graphicx}
\usepackage{dcolumn}
\usepackage{bm}
\usepackage{mathtools}
\usepackage{braket}
\usepackage{tikz}
\usetikzlibrary{arrows.meta,positioning,calc,decorations.pathmorphing}
\usepackage{quantikz}
\usepackage{xr-hyper}
\usepackage{hyperref}
\externaldocument[M-]{main}
\usepackage{amsthm}

% APS-style theorem formatting
\newtheoremstyle{aps}%
  {}{}%
  {}%
  {\parindent}%
  {\itshape}%
  {.}%
  { }%
  {}%

\theoremstyle{aps}
\newtheorem{theorem}{Theorem}
\newtheorem{lemma}[theorem]{Lemma}
\newtheorem{proposition}[theorem]{Proposition}
\newtheorem{corollary}[theorem]{Corollary}
\newtheorem{definition}[theorem]{Definition}
\newtheorem{example}[theorem]{Example}
\newtheorem{remark}[theorem]{Remark}

% Custom commands
\newcommand{\C}{\mathbb{C}}
\newcommand{\R}{\mathbb{R}}
\newcommand{\N}{\mathbb{N}}
\newcommand{\Z}{\mathbb{Z}}
\newcommand{\T}{\mathbb{T}}
\newcommand{\Sphere}{\mathbb{S}^2}
\newcommand{\CP}{\mathbb{C}\mathrm{P}^1}
\newcommand{\SU}{\mathrm{SU}}
\newcommand{\SO}{\mathrm{SO}}
\newcommand{\U}{\mathrm{U}}
\newcommand{\Tr}{\mathrm{Tr}}
\newcommand{\re}{\mathrm{Re}}
\newcommand{\im}{\mathrm{Im}}

\newcommand{\PauliI}{\sigma_0}
\newcommand{\PauliX}{\sigma_x}
\newcommand{\PauliY}{\sigma_y}
\newcommand{\PauliZ}{\sigma_z}

\newcommand{\rtilde}{\tilde{r}}
\newcommand{\TB}{\mathrm{TB}}
\newcommand{\SB}{\mathrm{SB}}


% Supplemental numbering
\renewcommand{\theequation}{S\arabic{equation}}
\renewcommand{\thesection}{S\Roman{section}}
\renewcommand{\thetheorem}{S\arabic{theorem}}

\begin{document}

\title{Supplemental Material: Quantum state preparation on the real line with stereographic quantum signal processing}

\thanks{This manuscript has been authored by UT-Battelle, LLC, under Contract No.\ DE-AC05-00OR22725 with the U.S.\ Department of Energy (DOE). The United States Government retains and the publisher, by accepting the article for publication, acknowledges that the United States Government retains a non-exclusive, paid-up, irrevocable, world-wide license to publish or reproduce the published form of this manuscript, or allow others to do so, for United States Government purposes. The Department of Energy will provide public access to these results of federally sponsored research in accordance with the DOE Public Access Plan.}

\author{Tyler Kharazi}
\affiliation{Department of Chemistry, University of California, Berkeley, California 94720, USA}

\author{Vicente Leyton-Ortega}
\email{leytonorteva@ornl.gov}
\affiliation{Quantum Computational Science Group, Quantum Information Science Section, Computational Sciences and Engineering Division, Oak Ridge National Laboratory, Oak Ridge, Tennessee 37831, USA}

\date{\today}

\maketitle

\section{Direct proof of Theorem~\ref{M-thm:real-qsp} of the main text}
\label{app:thm-real}

The proof in Sec.~\ref{M-subsec:signal} of the main text reduces the theorem to the
standard one.  Here we prove it directly in the variable $x$, by
induction on $d$.

\emph{Forward direction.}  For $d=0$,
$e^{i\phi_0\PauliZ}=\mathrm{diag}(e^{i\phi_0},e^{-i\phi_0})$ has
$P=e^{i\phi_0}$, $Q\equiv 0$, and $|P|^2+|Q|^2=1=(1+x^2)^0$, with
$\deg P=0$ and $\deg Q=-1$.  Assume the form of Eq.~\eqref{M-eq:real-qsp-form} of the main text
at degree $d-1$.  Multiplying by $e^{i\phi_d\PauliZ}O_x$ gives
\begin{equation}
	\begin{pmatrix}\widetilde P\\ \widetilde Q\end{pmatrix}
	=\begin{pmatrix} e^{i\phi_d}(xP-Q)\\ e^{-i\phi_d}(xQ+P)\end{pmatrix},
\end{equation}
so that $e^{i\phi_d\PauliZ}O_x U_{\bm\phi}=
\bigl(\begin{smallmatrix}\widetilde P&-\widetilde Q^*\\
\widetilde Q&\widetilde P^*\end{smallmatrix}\bigr)$.  Direct expansion
gives
$|\widetilde P|^2+|\widetilde Q|^2=|xP-Q|^2+|xQ+P|^2
=(1+x^2)(|P|^2+|Q|^2)=(1+x^2)^d$, the cross terms cancelling.  Since
$x$ multiplies $P$ on the diagonal and $Q$ on the off-diagonal,
$\deg\widetilde P\le d$ and $\deg\widetilde Q\le d-1$, and the
parities shift by one as required.

\emph{Backward direction.}  Let $P,Q$ satisfy (i)--(iii).  For
real $x$, $|P(x)|^2=P(x)P^*(x)$ with $P^*(x)=\overline{P(\bar x)}$,
so (i) states that the polynomial $PP^*+QQ^*-(1+x^2)^d$ vanishes on
$\R$.  A polynomial that vanishes on $\R$ is identically zero, so
$P(x)P^*(x)+Q(x)Q^*(x)=(1+x^2)^d$ holds for all complex $x$, and in
particular at $x=i$, where the right side vanishes.  There
$P^*(i)=\overline{P(-i)}=(-1)^d\overline{P(i)}$ and
$Q^*(i)=(-1)^{d-1}\overline{Q(i)}$ by the parities (ii), so the
identity at $x=i$ reads $(-1)^d\bigl(|P(i)|^2-|Q(i)|^2\bigr)=0$,
that is, $|P(i)|^2=|Q(i)|^2$.  The outermost factor
$e^{i\phi\PauliZ}O_x$ is removed by applying
$(1+x^2)\,O_x^{-1}e^{-i\phi\PauliZ}$, with
$(1+x^2)\,O_x^{-1}=\bigl(\begin{smallmatrix}x&1\\-1&x
\end{smallmatrix}\bigr)$, which gives
\begin{equation}
	\begin{pmatrix}\widetilde P\\ \widetilde Q\end{pmatrix}
	=\begin{pmatrix} x\,e^{-i\phi}P+e^{i\phi}Q\\
		-e^{-i\phi}P+x\,e^{i\phi}Q\end{pmatrix}.
\end{equation}
If $Q(i)\neq 0$, choose $\phi$ with $e^{2i\phi}=-iP(i)/Q(i)$, which
has unit modulus by $|P(i)|=|Q(i)|$.  Then
$\widetilde P(i)=e^{-i\phi}\bigl(iP(i)+e^{2i\phi}Q(i)\bigr)=0$, and
$\widetilde Q(i)=i\,\widetilde P(i)=0$.  By the parities (ii),
$\widetilde P(-i)=(-1)^{d+1}\widetilde P(i)$ and
$\widetilde Q(-i)=(-1)^{d}\widetilde Q(i)$, so both vanish also at
$x=-i$ and $1+x^2$ divides $\widetilde P$ and $\widetilde Q$.  Since
$|\widetilde P|^2+|\widetilde Q|^2=(1+x^2)^{d+1}$, the
quotients $\widehat
P=\widetilde P/(1+x^2)$, $\widehat Q=\widetilde Q/(1+x^2)$ satisfy
$|\widehat P|^2+|\widehat Q|^2=(1+x^2)^{d-1}$ with degrees $d-1$ and
$d-2$.  If instead $Q(i)=0$, then $P(i)=0$ as well, $1+x^2$
already divides both, and any $\phi$ works.  Recursing to $d=0$
yields the phases $\bm\phi$.\hfill$\square$

\section{Complex signals and the nonlinear Fourier transform}
\label{app:thm-gen}

Allowing a complex signal $z\in\C$ and arbitrary $\SU(2)$ rotations
$G_0,\dots,G_d$ in place of the phase rotations gives the
stereographic form of generalized QSP~\cite{MotlaghWiebe2023}.  The
signal operator is
\begin{equation}
	O_z=\begin{pmatrix} z & -1 \\ 1 & \bar z\end{pmatrix},
	\qquad \det O_z=1+|z|^2,
	\label{eq:Oz}
\end{equation}
which reduces to $O_x$ of Eq.~\eqref{M-eq:Ox} of the main text for real $z$.  The
entries of a product of such factors are polynomials in $z$ and
$\bar z$.  For such a polynomial, $P^*(z,\bar z)=
\overline{P(\bar z,z)}$ conjugates the coefficients and exchanges
$z$ and $\bar z$, so that $P^*=\overline{P}$ for every $z\in\C$.

\begin{proposition}[Generalized stereographic sequences]
	\label{thm:gen-qsp}
	For all $G_0,\dots,G_d\in\SU(2)$,
	\begin{equation}
		G_d\,O_z\,G_{d-1}\,O_z\cdots O_z\,G_0
		=\begin{pmatrix} P & -Q^*\\ Q & P^*\end{pmatrix}
		\quad\text{for all }z\in\C,
		\label{eq:gen-qsp-form}
	\end{equation}
	where $P$ and $Q$ are polynomials in $z$ and $\bar z$ such that
	(i) every monomial $z^j\bar z^k$ has $j+k\le d$, and (ii)
	$|P|^2+|Q|^2=(1+|z|^2)^d$ for all $z\in\C$.
\end{proposition}

\begin{proof}
	We induct on $d$.  For $d=0$ the product is $G_0\in\SU(2)$.
	Let $M_{d-1}$ have the stated form at degree $d-1$.  Both
	$O_z$ and $G_d$ have the form
	$\bigl(\begin{smallmatrix}\alpha&-\bar\beta\\
	\beta&\bar\alpha\end{smallmatrix}\bigr)$, which is closed
	under multiplication, so $G_dO_zM_{d-1}$ has the form of
	Eq.~\eqref{eq:gen-qsp-form}.  Its entries are those of
	$M_{d-1}$ multiplied by $z$, $\bar z$, or constants, which
	gives (i), and its determinant is
	$\det O_z\,\det M_{d-1}=(1+|z|^2)^d$, which is (ii).
\end{proof}

We do not characterize which pairs obeying (i) and (ii) arise in
this way.

\emph{Relation to the nonlinear Fourier transform.}  The $\SU(2)$
nonlinear Fourier transform (NLFT) of a sequence
$(\gamma_0,\dots,\gamma_{d-1})$ is the ordered product
\begin{equation}
	\widehat\gamma(w)=\prod_{k=0}^{d-1}
	\frac{1}{\sqrt{1+|\gamma_k|^2}}
	\begin{pmatrix} 1 & \gamma_k w^k\\ -\bar\gamma_k w^{-k} & 1
	\end{pmatrix},\qquad |w|=1,
	\label{eq:nlft}
\end{equation}
which is equivalent to generalized
QSP~\cite{AlexisMnatsakanyanThiele2024,AlexisBeckerEtAl2026,Laneve2025b}.
Each unnormalized factor at $w=1$ is
$\bigl(\begin{smallmatrix}1&\gamma\\-\bar\gamma&1\end{smallmatrix}\bigr)
=O_\gamma\,(i\PauliY)$ with $i\PauliY=\bigl(\begin{smallmatrix}0&1\\
-1&0\end{smallmatrix}\bigr)$, which explains why the same
normalization $\sqrt{1+|\gamma|^2}$ appears in both constructions.
We note that this identity says only that both matrices are
positive multiples of $\SU(2)$ matrices of the same shape, and that
$\gamma$ plays the role of a coefficient in the NLFT but of the
signal in $O_\gamma$.  The inverse problem of the NLFT is solved by
layer stripping~\cite{NiSarkarYingLin2025,AlexisBeckerEtAl2026}.

\bibliography{references,references_extra}

\end{document}
